\hwhat{Mean-field Glauber dynamics ties the gain $g$ read from equal-time fluctuations to how much slower the collective mode relaxes than one agent. Round 1 (70 units) found that in regime III the collective talk mode keeps more lag-1 memory than this predicts ($\Delta\rho_1$, post hoc). The open rival was a slow common drive, the scheduler. Round 2 re-indexed talk on each agent's call clock, removed the scheduler field, and added partition contrasts. \textbf{The excess is read-out coupling, not the scheduler.} It survives every removal ($+0.052$ median; \#51 $+0.042\,[0.014,0.070]$), lives inside rooms ($\Pi\approx15$), and equals what H67's read-out gain produces ($g_{1,\rm fit}/g_{lag}=0.96$). The $\times9.5$ nudge result is withdrawn. The holdout is not run.}

\hmodel{Models 02, 01 and 09. Talk spin $s_i(t)$ per minute on the call-based all-present window, or $Y_{ic}$ per receiving call $c$. Read-out coupling acts at the recipient's next call: messages posted before $t_c$ are read; messages posted during the call are in flight. A slow common drive acts on every agent at once, whether or not it reads. Each unit's null is the W0 world: independent talk at its real calls.}

\hmath{\vspace{-6pt}\begin{gather*}
\Delta\rho_1=\rho_c(1)-\rho_\perp(1)^{1-g_\chi},\qquad \Delta\rho_1^{c}=\Delta\rho_1-\langle\Delta\rho_1\rangle_{W0}\\
P(Y_{ic}{=}1)=a_{i,d,\kappa}+\textstyle\sum_h\beta_h R^h_{ic}+\beta_0P^0_{ic}+\beta_X X^1_{ic}+\rho\,Y_{i,c-1}\\
g_1=\bar m\,\bar r\,(\beta_1-\beta_X),\qquad \Omega=\textstyle\sum_k G(k)\big/\big[G(p)\sum_k S(k)/S(p)\big]
\end{gather*}\vspace{-8pt}}

\hobs{../figures/r2_obs.pdf}{Left: one point per regime-III unit (size $\propto$ trimmed calls). The W0-corrected collective memory before (V0) and after (V5) removing the call-window, day-edge, exogenous and activity-field structure. Points stay near the diagonal, so the removals change little. Right: the hop-1 gain that a call-skeleton simulator needs to reproduce each unit's minute memory, against H67's read-out gain. Bars are CIs; dotted lines mark $\times2$.}

\htelos{For a monitor: in regime III, read the collective talk memory as read-out coupling of size $g\approx0.11$--$0.14$ (push amplified $\times1.15$), not as a scheduler artefact. Report $\Delta\rho_1^c$ next to $g_{lag}$; they agree. For an operator: naming the recipient is the lever. A named message raises the reader's talk chance by 7.3 points at its next call and by 2.5 points one call later; an unnamed one by 0.2. A nudge's response is about what the target's own spontaneous wake predicts; it is not tenfold.}

\hbottom{The regime-III talk memory that mean-field Glauber misses is coupling delayed by one read-out call. It survives scheduler removal, respects the room and address gates, and matches H67's read-out gain. The tenfold nudge excess was an alignment artefact.}

\hresults{\scriptsize\setlength{\tabcolsep}{2pt}\begin{tabular}{@{}p{0.33\columnwidth}p{0.50\columnwidth}p{0.15\columnwidth}@{}}
\textbf{Prediction (dated)} & \textbf{Observed} & \textbf{Verdict}\\\hline
R1 synthetic (real skeletons) & $J_1^*$ only in coupling worlds; $g_1$ recovered +10\%; room partition fails only under a room-local drive & 3/5 pass, 2 marginal\\
Q-S1 survives removal (V5) & median $+0.052$; \#51 $+0.042\,[0.014,0.070]$ & supported\\
Q-S2 no tracking of $c_\times$ & Spearman $-0.10$ ($p$ 0.64) & supported\\
Q-P1 room gate (minute grid) & same 0.0165, cross 0.0011 [$-$0.003, 0.005] & supported\\
Q-P2 address gate (call clock) & named $\times10$; placebos negative & mixed\\
Q-R1a self memory $>0$ & $-0.048$ (refractory) & failed\\
Q-R1b not instantaneous & in flight $-0.015$ vs read $+0.073$ (named) & supported\\
Q-R1c kernel $>$ one call & named 0.025 (ratio 0.34); unnamed 0 & mixed\\
Q-R2 $g_{1,\rm fit}\approx g_{lag}$ & ratio 0.96 (median, 24 units) & supported\\
Q-R4a nudge $\Omega\le2$ & 1.29; $A_{30}$ 0.44 [$-$0.41, 1.23] (post hoc) & supported\\
Q-R4b human $\Omega\in[0.5,2]$ & talk per call 2.53 [1.63, 3.44] (post hoc) & mixed\\
\end{tabular}}

\hobsb{../figures/r2_kernel.pdf}{Left: regime-III pooled talk response per message by window on the call clock. Messages a call cannot read (in flight, other rooms) carry a small negative association. Read messages carry a positive one, almost all from messages that name the recipient, which keep acting one call later. Right: \#51 nudge activity response aligned on the target's receiving call, with the matched spontaneous regression (Onsager) scaled to its peak.}

\hcaveats{$\Delta\rho_1$ was chosen after round-1 data (A2). The round-2 tests were pre-registered, with estimator fixes found in the synthetic (B1). The room partition cannot exclude a room-local drive; the call-clock gates do. The in-flight placebo is negative in regime III, so the in-flight-gated gain (0.23) overstates; the room-gated one (0.11) is used. R4 normalization changed after a first run (B2, post hoc), and the nudge activity CI is wide. Regime I stays untestable on this design.}

\hnext{Holdout on \#43 (V5 $\Delta\rho_1^c$, room partition, $g_{1,\rm fit}/g_{lag}$) via a re-frozen confirm script. Fit the full per-call kernel of named messages against H40's decay. Identify the negative in-flight field.}
